\section{Module-Oriented Database Design}\label{sec:module-design}

The following views show the tables each product module reads or writes. Shared entities may appear in more than one view because module ownership and relational ownership are not identical.

\subsection{Module 1 - Medication Intake}

Module 1 captures medication information from text, image, or document input. A user creates \db{MEDICATION_INTAKES}; each submission may contain \db{INTAKE_DOCUMENTS}. One or more \db{EXTRACTION_RUNS} consume those documents through \db{EXTRACTION_RUN_INPUTS}. Parsed rows become \db{MEDICATION_ENTRIES}, optionally grouped by \db{PRESCRIPTIONS}, and the patient records confirmation or correction in \db{ENTRY_CONFIRMATIONS}.

\portraitdiagram{module-01-medication-intake.png}{Module 1 - Medication Intake ERD}{fig:module-intake}

\begin{longtable}{L{4.5cm}L{10.8cm}}
\caption{Module 1 tables and responsibilities}\label{tab:module1-tables}\\
\toprule
\textbf{Table} & \textbf{Responsibility}\\
\midrule
\endfirsthead
\toprule
\textbf{Table} & \textbf{Responsibility}\\
\midrule
\endhead
\db{USERS} & Identifies the submitter without duplicating account data in the intake domain.\\
\db{MEDICATION_INTAKES} & Represents one intake session and records input method, owner, status, and submission context.\\
\db{INTAKE_DOCUMENTS} & Stores metadata and secure object references for images or documents; binary content should reside in object storage.\\
\db{EXTRACTION_RUNS} & Records an OCR/parser execution, model or rule version, timestamps, and outcome.\\
\db{EXTRACTION_RUN_INPUTS} & Many-to-many bridge between extraction runs and the documents they process.\\
\db{PRESCRIPTIONS} & Groups medication entries that belong to the same patient prescription or source document.\\
\db{MEDICATION_ENTRIES} & Holds raw text, extracted fields, status, and the patient-facing medication row.\\
\db{ENTRY_CONFIRMATIONS} & Keeps the actor, decision, correction, and time of confirmation instead of overwriting extraction history.\\
\bottomrule
\end{longtable}

\subsection{Module 2 - Medication Identification and Reference}

Module 2 resolves a confirmed medication entry against one released medication catalog. \db{ENTRY_IDENTIFICATIONS} records an identification attempt and selected result; \db{IDENTIFICATION_CANDIDATES} retains alternatives and scores. The selected \db{MEDICATIONS} row expands through components, ingredients, strengths, indications, and conditions.

\landscapediagram{module-02-medication-identification-reference.png}{Module 2 - Medication Identification and Reference ERD}{fig:module-identification}

\begin{longtable}{L{4.5cm}L{10.8cm}}
\caption{Module 2 principal tables and responsibilities}\label{tab:module2-tables}\\
\toprule
\textbf{Table} & \textbf{Responsibility}\\
\midrule
\endfirsthead
\toprule
\textbf{Table} & \textbf{Responsibility}\\
\midrule
\endhead
\db{DATASET_RELEASES} & Selects the approved reference-data release used by the identification attempt.\\
\db{ENTRY_IDENTIFICATIONS} & Records input entry, release, method, status, confidence, selected medication, and reviewer state.\\
\db{IDENTIFICATION_CANDIDATES} & Preserves ranked candidates so ambiguity can be reviewed and evaluated.\\
\db{MEDICATIONS} & Versioned product identity, including name, registration, manufacturer, dosage form, route, and verification status.\\
\db{INGREDIENTS} & Versioned ingredient identity with source-local and accepted normalized names.\\
\db{MEDICATION_COMPONENTS} & Connects a medication to a component and, when accepted, to an ingredient identity.\\
\db{STRENGTH_ASSERTIONS} & Stores raw and normalized strength assertions separately from component identity.\\
\db{MEDICATION_INDICATIONS} & Stores source indication statements and approval state.\\
\db{CONDITIONS} & Stores reviewed condition concepts.\\
\db{INDICATION_CONDITIONS} & Many-to-many link between approved indication statements and normalized conditions.\\
\db{MEDICATION_PRIORITIES} & Stores collection/review priority; it must not be interpreted as clinical recommendation.\\
\bottomrule
\end{longtable}

\subsection{Module 3 - Medication Safety Assessment}

Module 3 freezes the patient's current medication list in \db{MEDICATION_LIST_SNAPSHOTS}. \db{SNAPSHOT_ITEMS} reference the exact entries and identifications included. A \db{SAFETY_ASSESSMENTS} row binds that snapshot to a dataset release. Pairwise results are retained in \db{ASSESSMENT_PAIR_RESULTS}; user-facing findings are stored in \db{SAFETY_FINDINGS} and linked to affected items through \db{FINDING_ITEMS}.

\landscapediagram{module-03-medication-safety-assessment.png}{Module 3 - Medication Safety Assessment ERD}{fig:module-safety}

\begin{longtable}{L{4.5cm}L{10.8cm}}
\caption{Module 3 principal tables and responsibilities}\label{tab:module3-tables}\\
\toprule
\textbf{Table} & \textbf{Responsibility}\\
\midrule
\endfirsthead
\toprule
\textbf{Table} & \textbf{Responsibility}\\
\midrule
\endhead
\db{MEDICATION_LIST_SNAPSHOTS} & Immutable assessment-time list owned by a patient.\\
\db{SNAPSHOT_ITEMS} & Connects snapshot rows to medication entries and accepted identification results.\\
\db{SAFETY_ASSESSMENTS} & Assessment execution, dataset release, lifecycle status, and timestamps.\\
\db{ASSESSMENT_PAIR_RESULTS} & Canonical unordered pair evaluation with matched interaction and outcome.\\
\db{SAFETY_FINDINGS} & Patient/HCP-facing finding with severity, explanation, and referenced interaction.\\
\db{FINDING_ITEMS} & Links one finding to the medication items involved.\\
\db{SAFETY_DRUG_ENTITIES} & Reviewed endpoint identifiers from the interaction namespace.\\
\db{COMPONENT_SAFETY_MAPPINGS} & Reviewed component-to-safety-entity mapping; candidate mappings are non-operative.\\
\db{DRUG_INTERACTIONS} & Canonical interaction record independent of source-specific participant order.\\
\db{DRUG_INTERACTION_PARTICIPANTS} & Interaction endpoints; the physical model must enforce exactly two distinct participants per pair.\\
\db{INTERACTION_EVIDENCE} & Connects interactions to one or more evidence sources and preserves disagreement.\\
\bottomrule
\end{longtable}

\subsection{Module 4 - HCP Review}

Module 4 creates an explicit request for professional interpretation of a safety assessment. The request records patient, assessment, care relationship, assignee, priority, and status. \db{HCP_REVIEWS} retains the professional decision. A linked conversation permits the patient and assigned HCP to communicate directly.

\portraitdiagram{module-04-hcp-review.png}{Module 4 - HCP Review ERD}{fig:module-hcp-review}

\begin{longtable}{L{4.5cm}L{10.8cm}}
\caption{Module 4 tables and responsibilities}\label{tab:module4-tables}\\
\toprule
\textbf{Table} & \textbf{Responsibility}\\
\midrule
\endfirsthead
\toprule
\textbf{Table} & \textbf{Responsibility}\\
\midrule
\endhead
\db{CARE_RELATIONSHIPS} & Establishes the professional's authorized patient scope.\\
\db{HCP_REVIEW_REQUESTS} & Case assignment and workflow state for review of one assessment.\\
\db{HCP_REVIEWS} & Professional interpretation, decision, recommendations, and completion metadata.\\
\db{CONVERSATIONS} & Case-bound communication channel, not a general public chat.\\
\db{CONVERSATION_PARTICIPANTS} & Explicit participant membership and access boundary.\\
\db{MESSAGES} & Immutable message content and sender metadata within the authorized conversation.\\
\bottomrule
\end{longtable}

\subsection{Module 5 - Medication Reminders}

Module 5 schedules reminders only from a regimen associated with patient entries and, where required, HCP review. Eligibility checks prevent an unconfirmed or clinically blocked item from producing reminder events. Adherence is recorded separately from delivery events.

\portraitdiagram{module-05-medication-reminders.png}{Module 5 - Medication Reminders ERD}{fig:module-reminders}

\begin{longtable}{L{4.5cm}L{10.8cm}}
\caption{Module 5 tables and responsibilities}\label{tab:module5-tables}\\
\toprule
\textbf{Table} & \textbf{Responsibility}\\
\midrule
\endfirsthead
\toprule
\textbf{Table} & \textbf{Responsibility}\\
\midrule
\endhead
\db{MEDICATION_REGIMENS} & Patient regimen header and lifecycle state.\\
\db{REGIMEN_ITEMS} & Instructions and date range for one medication entry within a regimen.\\
\db{REGIMEN_REVIEW_LINKS} & Records which HCP review supports a regimen.\\
\db{REMINDER_SCHEDULES} & Frequency, timing, timezone, start/end, and schedule status.\\
\db{SCHEDULE_ELIGIBILITY_CHECKS} & Audit of why a schedule is eligible, blocked, or requires review.\\
\db{REMINDER_EVENTS} & Concrete planned/sent/failed reminder occurrence; supports idempotent delivery.\\
\db{ADHERENCE_RECORDS} & Patient or authorized-user response such as taken, skipped, or missed.\\
\db{NOTIFICATIONS} & Delivery envelope for reminders and other workflow notifications.\\
\bottomrule
\end{longtable}

